\chapter{自由落体与抛体运动}\label{ch:g12:projectile-motion}

花园水管喷出的每一股水流，在空中都画出同样的曲线。这一曲线也出现在
飞向篮筐的篮球、跳远运动员、炮弹上：凡是只受\omterm{def:g10:universal-gravitation:weight}{重力}作用的物体，
都沿这一曲线运动——即抛物线——本章由\omterm{def:g10:inertia:force}{牛顿}第二定律出发，用两行微积分推出它，
再从其方程读出最高点、\omterm{def:g12:projectile-motion:range}{飞行时间}和\omterm{def:g12:projectile-motion:range}{射程}，
然后让空气重新回到模型中，看它
会破坏什么。

\section{自由落体}

\begin{definition}[自由落体]\label{def:g12:projectile-motion:freefall}
物体处于\emph{自由落体}\index{自由落体}状态，是指它所受的唯一的力
为其\omterm{def:g10:universal-gravitation:weight}{重力} $\vect P = m\vect g$，无支持力，无绳索拉力，空气阻力
可忽略。被释放的石块，或离开手的球，在很好的近似下
都处于自由落体状态——“落体”一词也包括向上运动。
\end{definition}

\begin{theorem}[伽利略自由落体定律]\label{thm:g12:projectile-motion:galileo}
\index{伽利略自由落体定律}
在\omterm{def:g12:projectile-motion:freefall}{自由落体}中，所有物体的加速度相同，与质量无关：
\[
\vect a = \vect g, \qquad g = \qty{9.81}{m/s^2},
\]
方向竖直向下：这里的场强 \qty{9.81}{N/kg} 即
在\cref{ch:g10:universal-gravitation}中的场强，只是写法不同，$\unit{N/kg} = \unit{m/s^2}$。
\end{theorem}

\begin{proof}
由\omterm{def:g10:inertia:force}{牛顿}第二定律（\cref{ch:g12:newtons-laws}）
并取\omterm{def:g10:universal-gravitation:weight}{重力}为唯一的力：$m\vect a = m\vect g$，质量相消。锤子与羽毛
的差别仅来自空气。
\end{proof}

\begin{proposition}[从静止开始的竖直下落]\label{prop:g12:projectile-motion:vertical}
物体从静止释放，下落深度 $h$ 自释放点向下计，
则处于\omterm{def:g12:projectile-motion:freefall}{自由落体}中的物体满足
\[
a = g, \qquad v = g t, \qquad h = \tfrac12\, g t^2 .
\]
\end{proposition}

\begin{proof}
$v(t)$ 是常量 $a = g$ 的原函数，且在 $t = 0$ 时为零；
$h(t)$ 是 $v = gt$ 的原函数，且在 $t = 0$ 时为零（见\cref{ch:g12:kinematics-2d}）。
\end{proof}

\begin{example}[数量级]\label{ex:g12:projectile-motion:orders}
\omterm{def:g12:projectile-motion:freefall}{自由落体}经过 $1$、$2$、$3$ 秒后：$v = 9.8$、$19.6$、\qty{29.4}{m/s}
（$35$、$71$、\qty{106}{km/h}），$h = 4.9$、$19.6$、\qty{44.1}{m}。从
\qty{30}{m} 高的屋顶下落：$t = \sqrt{2h/g} \approx \qty{2.5}{s}$，落地速率为
\qty{24}{m/s} $\approx \qty{87}{km/h}$。时间随 $\sqrt h$ 增长——而
经过几秒后，空气就不能再忽略（见上一节）。
\end{example}

\begin{example}[竖直上抛]\label{ex:g12:projectile-motion:throw}
以 $v_0$ 竖直向上抛出，取 $Oy$ 轴向上，原点在手处。同样的两次
求原函数给出 $v = v_0 - gt$ 与 $y = v_0 t - \tfrac12 gt^2$。最高点
即 $v = 0$ 处，出现在 $t = v_0/g$，高度 $h = v_0^2/2g$——正是
将\cref{ch:g11:mechanical-energy}中的\omterm{def:g10:energy-conservation:energy}{能量}结论 $v_0 = \sqrt{2gh}$ 倒过来使用；
若\emph{向下}抛出，则 $v_0$ 变号。取 \qty{15}{m/s}：最高点为
\qty{11.5}{m}，出现在 \qty{1.53}{s}，\qty{3.06}{s} 时回到出发点，速率为 \qty{15}{m/s}。
\end{example}

\section{抛体：同一定律，两次求原函数}

\begin{definition}[抛体]\label{def:g12:projectile-motion:projectile}
\emph{抛体}\index{抛体}是指以速度
$\vect v_0$ 抛出、此后仅受\omterm{def:g10:universal-gravitation:weight}{重力}作用的物体。以抛出点为原点建立坐标轴：$Ox$
轴水平，位于 $\vect v_0$ 所在的竖直平面内，$Oy$ 轴竖直向上；
\emph{抛射角}\index{抛射角} $\alpha$ 为 $\vect v_0$ 与水平方向的夹角，
它决定了分量 $(v_0\cos\alpha,\; v_0\sin\alpha)$。
\end{definition}

\begin{theorem}[运动方程]\label{thm:g12:projectile-motion:equations}
以速率 $v_0$、夹角 $\alpha$ 从原点抛出的\omterm{def:g12:projectile-motion:projectile}{抛体}，
在任一时刻都有 $\vect a = (0,\,-g)$，因此
\[
\vect v(t) = \bigl(v_0\cos\alpha,\; v_0\sin\alpha - g t\bigr),
\qquad
\begin{cases}
x(t) = v_0\cos\alpha\; t,\\[2pt]
y(t) = v_0\sin\alpha\; t - \tfrac12\, g t^2 .
\end{cases}
\]
\end{theorem}

\begin{proof}
由\omterm{def:g10:inertia:force}{牛顿}第二定律得 $\vect a = \vect g = (0,-g)$。逐个坐标轴求原函数：
$\vect v(t)$ 由 $\vect a$ 及初值
$(v_0\cos\alpha,\, v_0\sin\alpha)$ 求出，再由 $\vect v$ 及 $(0,0)$ 求出位置。
\end{proof}

\begin{proposition}[两种运动的独立性]\label{prop:g12:projectile-motion:independence}
水平运动是匀速的——$v_x$ 始终不变——而
竖直运动恰好是上一节的竖直上抛：两个方程都不涉及对方的坐标。
特别地，一个水平射出的球，
与一个在同一时刻、从同一高度释放的球，
始终处于同一水平线上——它们同时落地。
\end{proposition}

\begin{proof}
由\cref{thm:g12:projectile-motion:equations}可知，$x(t)$ 只与
$v_0\cos\alpha$ 有关，$y(t)$ 只与 $v_0\sin\alpha$ 及 $g$ 有关；两球都有
$v_y(0) = 0$，因此两者都满足 $h = \tfrac12 gt^2$。
\end{proof}

\begin{omfigure}
\begin{tikzpicture}[scale=0.9, font=\small]
  \draw[omExo] (-0.7,0) -- (3.5,0);
  \draw[dashed, omExo] (0,3) -- (0,0.75);
  \draw[thick, omThm] plot[domain=0:2.7, samples=50] (\x,{3-0.3086*\x*\x});
  \foreach \t/\y in {0/3.0, 1/2.75, 2/2.0, 3/0.75}
    { \fill[omDef] (0,\y) circle (2.5pt); \fill[omThm] ({0.9*\t},\y) circle (2.5pt); }
  \foreach \t/\y in {1/2.75, 2/2.0, 3/0.75}
    \draw[dashed, omExo] (0.06,\y) -- ({0.9*\t-0.08},\y);
  \draw[-Stealth, omThm] (0,3) -- (0.8,3) node[right] {$\vect v_0$};
  \node[left, omDef] at (-0.15,2.75) {$t_1$};
  \node[left, omDef] at (-0.15,2.0) {$t_2$}; \node[left, omDef] at (-0.15,0.75) {$t_3$};
\end{tikzpicture}

{\small 一球自由下落，一球同时水平射出：在
每个时刻 $t_1$、$t_2$、$t_3$ 两球高度相同——速率不造成任何\omterm{def:g12:mechanical-waves:delay}{延迟}。}
\end{omfigure}

\begin{remark}[下落的猴子]\label{rem:g12:projectile-motion:monkey}
动物园饲养员用飞镖\emph{直接瞄准}一只猴子，猴子
在射击时松开了树枝。这样无济于事：飞镖与猴子各自下落 $\tfrac12 gt^2$，
都低于各自无\omterm{def:g10:universal-gravitation:weight}{重力}时的位置——即瞄准线与树枝——因此
无论飞镖速率如何，二者恰好相遇
（见\cref{exo:g12:projectile-motion:15}）。
\end{remark}

\section{完整地看抛物线}

\begin{theorem}[轨迹]\label{thm:g12:projectile-motion:trajectory}
消去 $t$：由 $x(t)$ 和 $y(t)$ 可知，\omterm{def:g12:projectile-motion:projectile}{抛体}
（$0 \leq \alpha < \ang{90}$）的\omterm{def:g10:relative-motion:trajectory}{轨迹}是开口向下的抛物线
\[
y = x\tan\alpha \;-\; \frac{g}{2 v_0^2 \cos^2\alpha}\, x^2 .
\]
\end{theorem}

\begin{proof}
将 $t = x/(v_0\cos\alpha)$ 代入
\cref{thm:g12:projectile-motion:equations}中的 $y(t)$。
\end{proof}

\begin{omfigure}
\begin{tikzpicture}[scale=0.9, font=\small]
  \draw[omExo] (-0.4,0) -- (8.7,0);
  \draw[thick] plot[domain=0:8, samples=60] (\x,{1.08*\x-0.135*\x*\x});
  \fill (0,0) circle (1.5pt);
  \draw[-Stealth, omThm, thick] (0,0) -- (0.75,0.81) node[above left] {$\vect v_0$};
  \draw (0.7,0) arc (0:47:0.7); \node at (1.0,0.26) {$\alpha$};
  \draw[-Stealth, omDef] (2,1.62) -- ++(1.0,0);
  \draw[-Stealth, omProp] (2,1.62) -- ++(0,0.54) node[left] {$\vect v_y$};
  \draw[-Stealth, omThm, thick] (2,1.62) -- ++(1.0,0.54) node[above] {$\vect v$};
  \draw[-Stealth, omThm, thick] (4,2.16) -- ++(1.0,0) node[above] {$\vect v$};
  \draw[-Stealth, omDef] (6,1.62) -- ++(1.0,0) node[above right] {$\vect v_x$};
  \draw[-Stealth, omProp] (6,1.62) -- ++(0,-0.54);
  \draw[-Stealth, omThm, thick] (6,1.62) -- ++(1.0,-0.54) node[right] {$\vect v$};
  \draw[-Stealth, thick] (8.3,2.2) -- ++(0,-0.7) node[below] {$\vect g$};
\end{tikzpicture}

{\small 飞行过程中：$\vect v_x$ 始终不变；
$\vect v_y$ 逐渐减小，在最高点为零，随后反向——箭头倾斜，
抛物线也随之变化。}
\end{omfigure}

\begin{definition}[射程]\label{def:g12:projectile-motion:range}
对于平地上的\omterm{def:g12:projectile-motion:projectile}{抛体}，\emph{射程}\index{射程} $R$ 是
\omterm{def:g12:projectile-motion:projectile}{抛体}回到发射高度时所经过的水平距离；
\emph{飞行时间}\index{飞行时间} $t_f$ 即该段行程的时长。
\end{definition}

\begin{proposition}[最高点、飞行时间与射程]\label{prop:g12:projectile-motion:range}
在平地上，$0 < \alpha < \ang{90}$ 时：
\[
t_f = \frac{2 v_0 \sin\alpha}{g}, \qquad
h_{\text{最高点}} = \frac{(v_0\sin\alpha)^2}{2g}
\ \ \text{ 到达于 } t_f/2, \qquad
R = \frac{v_0^2 \sin(2\alpha)}{g}.
\]
\omterm{def:g12:projectile-motion:range}{射程}在 $\alpha = \ang{45}$ 时最大，此时 $R = v_0^2/g$；
互余角 $\alpha$ 与 $\ang{90} - \alpha$ 对应相同的\omterm{def:g12:projectile-motion:range}{射程}。
\end{proposition}

\begin{proof}
令 $y = 0$，得 $t\,(v_0\sin\alpha - \tfrac12 gt) = 0$：解得发射时刻与 $t_f$。
最高点 $v_y = 0$ 位于 $t = v_0\sin\alpha/g = t_f/2$ 处，代入
$y(t)$ 得最高点高度。于是 $R = x(t_f) = 2v_0^2\sin\alpha\cos\alpha/g =
v_0^2\sin(2\alpha)/g$。正弦在 $2\alpha = \ang{90}$ 时最大，且
$\sin\bigl(2(\ang{90}-\alpha)\bigr) = \sin(\ang{180} - 2\alpha) =
\sin(2\alpha)$。
\end{proof}

\begin{omfigure}
\begin{tikzpicture}
\begin{axis}[omaxis, width=8.8cm, height=5.4cm,
    xlabel={$x$ (\unit{m})}, ylabel={$y$ (\unit{m})},
    xmin=0, xmax=10.8, ymin=0, ymax=5.4,
    legend pos=north east, legend cell align=left]
  \addplot[omDef, thick, domain=0:5.10] {x*tan(15) - 0.04905*x^2/cos(15)^2};
  \addplot[omProp, thick, domain=0:8.83] {x*tan(30) - 0.04905*x^2/cos(30)^2};
  \addplot[omThm, very thick, domain=0:10.19] {x*tan(45) - 0.04905*x^2/cos(45)^2};
  \addplot[omProp, thick, dashed, domain=0:8.83] {x*tan(60) - 0.04905*x^2/cos(60)^2};
  \addplot[omDef, thick, dashed, domain=0:5.10] {x*tan(75) - 0.04905*x^2/cos(75)^2};
  \legend{\ang{15}, \ang{30}, \ang{45}, \ang{60}, \ang{75}}
\end{axis}
\end{tikzpicture}

{\small 发射速率相同，$v_0 = \qty{10}{m/s}$，五个角度：\ang{45}
投得最远；两组角度（\ang{15}、\ang{75}）与（\ang{30}、\ang{60}）
落点相同——一组低平而快，一组高而慢。}
\end{omfigure}

\begin{example}[球门球]\label{ex:g12:projectile-motion:kick}
门将以 $v_0 = \qty{25}{m/s}$，$\alpha = \ang{45}$ 开球：
$R = 25^2/9.81 \approx \qty{64}{m}$；实际解围的落点更近
\qty{55}{m} 左右——空气阻力使\omterm{def:g12:projectile-motion:range}{射程}缩短。
\end{example}

\begin{method}[抛体问题的梳理]\label{met:g12:projectile-motion:method}
\begin{enumerate}
  \item 以发射点为原点建立坐标轴：$Ox$ 水平，$Oy$ 竖直向上；
        写出 $\vect v_0$ 的分量。
  \item 检验是否为\omterm{def:g12:projectile-motion:freefall}{自由落体}（只受\omterm{def:g10:universal-gravitation:weight}{重力}）：则 $\vect a = (0,-g)$；
        两次积分，每次都代入初始条件。
  \item 问\emph{位置}时：消去 $t$（得\omterm{def:g10:relative-motion:trajectory}{轨迹}）；问
        \emph{时刻}时：保留 $t$。
  \item 对照关系：最高点 $\Leftrightarrow v_y = 0$；发射高度
        $\Leftrightarrow y = 0$；墙位于 $d \Leftrightarrow x = d$ 处；地面
        低于发射点 $\ell$ 处 $\Leftrightarrow y = -\ell$。
  \item 检验：飞行关于最高点对称；$R \leq v_0^2/g$；空气阻力
        只会使\omterm{def:g12:projectile-motion:range}{射程}变短。
\end{enumerate}
\end{method}

\section{空气的影响}

\begin{definition}[终端速度]\label{def:g12:projectile-motion:terminal}
空气阻力与速度方向相反，且随速率增大而增大。
因此，下落物体只加速到阻力与\omterm{def:g10:universal-gravitation:weight}{重力}平衡为止；此后
$\vect a = \vect 0$，它以恒定的\emph{终端
速度}\index{终端速度}下落。
\end{definition}

\begin{example}[终端速度的数量级]\label{ex:g12:projectile-motion:airorders}
全速下落的跳伞者：约 \qty{55}{m/s}（\qty{200}{km/h}）。一滴雨滴：
约 \qty{9}{m/s}——若在真空中下落 \qty{2000}{m}，则其速率将达到
$\sqrt{2gh} \approx \qty{2.0e2}{m/s}$。羽毛约为 \qty{1}{m/s}，
在几厘米内即达到——羽毛缓慢下落的秘诀在于：它整个
下落过程都\emph{处于}\omterm{def:g12:projectile-motion:terminal}{终端速度}；锤子则从不接近自身的\omterm{def:g12:projectile-motion:terminal}{终端速度}。
\end{example}

\begin{remark}[真实弹道]\label{rem:g12:projectile-motion:drag}
阻力在整个飞行过程中始终与速度方向相反：\omterm{def:g12:projectile-motion:range}{射程}变短，
最高点降低，对称性被破坏——下降段比上升段更陡，
远距离射击明显短于抛物线的预测。严格而定量的处理
需要求解一个微分方程，这在本科一年级分册中完成；
本学年，对于密度较大、速率适中的\omterm{def:g12:projectile-motion:projectile}{抛体}，抛物线是合适的模型——
如速率为 \qty{10}{m/s} 的篮球，而非
速率为 \qty{70}{m/s} 的高尔夫球。
\end{remark}

\begin{omfigure}
\begin{tikzpicture}
\begin{axis}[omaxis, width=8.6cm, height=4.8cm,
    xlabel={$x$}, ylabel={$y$}, xtick=\empty, ytick=\empty,
    xmin=0, xmax=8.4, ymin=0, ymax=2.1,
    legend pos=north east, legend cell align=left]
  \addplot[omDef, thick, dashed, domain=0:8] {0.9*x - 0.1125*x^2};
  \addplot[omThm, thick, smooth] coordinates
    {(0,0) (0.8,0.62) (1.6,1.08) (2.4,1.38) (3.2,1.5) (4,1.42)
     (4.7,1.15) (5.3,0.7) (5.8,0.15) (5.95,0)};
  \legend{抛物线（无空气）, 有空气阻力}
\end{axis}
\end{tikzpicture}

{\small 同一次发射，分别在有空气与无空气时：阻力缩短\omterm{def:g12:projectile-motion:range}{射程}、降低最高点，
并使\omterm{def:g12:projectile-motion:projectile}{抛体}落地时比起飞时更陡——抛物线的
对称性首先被破坏。}
\end{omfigure}

\section{练习}

\begin{exercise}[$\star$]\label{exo:g12:projectile-motion:1}
一块石头从桥上落下，经 \qty{2.0}{s} 后落入水中：求桥的高度，
以及落水时的速率（以 \unit{m/s} 和 \unit{km/h} 表示）？
\end{exercise}

\begin{exercise}[$\star$]\label{exo:g12:projectile-motion:2}
在月球上，宇航员同时释放一把锤子和一根羽毛：它们
同时落地。为何会这样？羽毛在地球上为何明显落后？
\end{exercise}

\begin{exercise}[$\star$]\label{exo:g12:projectile-motion:3}
忽略空气阻力，计算从 \qty{1.0}{m}，
\qty{10}{m}、\qty{100}{m} 高处下落的时间与落地速率。时间随高度如何增长？
\end{exercise}

\begin{exercise}[$\star$]\label{exo:g12:projectile-motion:4}
一球以 \qty{15}{m/s} 竖直向上抛出：求最高点的高度与时间，
回到手中所需的总时间，以及返回时的速率？
\end{exercise}

\begin{exercise}[$\star$]\label{exo:g12:projectile-motion:5}
一颗弹珠从高 \qty{0.80}{m} 的桌面以 \qty{2.5}{m/s} 的速率滚落。求
下落时间及落点与桌边的水平距离。
下落时间是否与弹珠的速率有关？
\end{exercise}

\begin{exercise}[$\star\star$]\label{exo:g12:projectile-motion:6}
由课程中的五个发射角图：哪些成对的发射角\omterm{def:g12:projectile-motion:range}{射程}相同，
哪个角度\omterm{def:g12:projectile-motion:range}{射程}最远？对于低于最大\omterm{def:g12:projectile-motion:range}{射程}的某一\omterm{def:g12:projectile-motion:range}{射程}，
两个可能的角度中哪一个\omterm{def:g12:projectile-motion:range}{飞行时间}更长？用 $t_f$ 说明理由。
\end{exercise}

\begin{exercise}[$\star\star$]\label{exo:g12:projectile-motion:7}
从 \qty{20}{m} 高的阳台，一球竖直\emph{向下}抛出，
速率为 \qty{5.0}{m/s}。求落地时间及落地速率，并
用\omterm{def:g10:energy-conservation:energy}{能量}守恒（见\cref{ch:g11:mechanical-energy}）检验后者。
\end{exercise}

\begin{exercise}[$\star\star$]\label{exo:g12:projectile-motion:8}
一\omterm{def:g12:projectile-motion:projectile}{抛体}以 \qty{40}{m/s} 的速率从水平地面抛出，$\alpha = \ang{30}$：
求\omterm{def:g12:projectile-motion:range}{飞行时间}、最高点高度与\omterm{def:g12:projectile-motion:range}{射程}？
\end{exercise}

\begin{exercise}[$\star\star$]\label{exo:g12:projectile-motion:9}
将跳远运动员视为\omterm{def:g12:projectile-motion:projectile}{抛体}：起跳速率为 \qty{10.5}{m/s}，
$\alpha = \ang{20}$。求\omterm{def:g12:projectile-motion:range}{射程}；跳远冠军的成绩接近 \qty{9}{m}，
——试指出质点模型遗漏的两个因素。
\end{exercise}

\begin{exercise}[$\star\star$]\label{exo:g12:projectile-motion:10}
在离平地 \qty{1.2}{m} 高处，一颗子弹由静止释放，另一颗相同的子弹
以 \qty{300}{m/s} 的速率水平射出。哪一颗先落地？计算
水平射出的子弹的下落时间及落点距离。
考查的是哪一原理？
\end{exercise}

\begin{exercise}[$\star\star$]\label{exo:g12:projectile-motion:11}
用代数方法证明 $\alpha$ 与 $\ang{90} - \alpha$ 的\omterm{def:g12:projectile-motion:range}{射程}相同，
再求使初速率为 \qty{25}{m/s} 的\omterm{def:g12:projectile-motion:projectile}{抛体}落在 \qty{50}{m}
远处的两个发射角。
\end{exercise}

\begin{exercise}[$\star\star\star$]\label{exo:g12:projectile-motion:12}
水从花园软管以 \qty{8.0}{m/s} 的速率射出。求水平方向的最大\omterm{def:g12:projectile-motion:range}{射程}、
竖直方向的最大高度，以及能浇灌
距离 \qty{5.0}{m} 处花坛的两个喷口角度。
\end{exercise}

\begin{exercise}[$\star\star\star$]\label{exo:g12:projectile-motion:13}
一个球从地面以 \qty{18}{m/s} 的速率踢出，\omterm{def:g12:projectile-motion:projectile}{抛射角} $\alpha = \ang{35}$，
射向距离 \qty{15}{m} 处、高 \qty{4.0}{m} 的篱笆。利用\omterm{def:g10:relative-motion:trajectory}{轨迹}方程
判断球能否越过篱笆；若能，求越过的高度。
\end{exercise}

\begin{exercise}[$\star\star\star$]\label{exo:g12:projectile-motion:14}
一名 \qty{80}{kg} 的跳伞者身体平展，以终端速率 \qty{55}{m/s} 下落。
(a) 求空气阻力？(b) 画出起跳后 $v(t)$ 的图像形状。(c) 一颗雨滴从
\qty{2.0}{km} 高处落下，真空中的落地速率与实际
落地速率 \qty{9}{m/s} 相比如何？
\end{exercise}

\begin{exercise}[$\star\star\star$]\label{exo:g12:projectile-motion:15}
下落的猴子问题，证明：一支飞镖从原点出发，恰好瞄准
距离为 $D$、高度为 $H$（$\tan\alpha = H/D$）的猴子，猴子在
$t = 0$ 时开始下落。证明：飞镖到达 $x = D$ 时，
其高度与猴子的高度相同——与 $v_0$ 无关。
\end{exercise}

\section{问题：篮球绝杀}

\begin{problem}[{周末问题——篮球绝杀：罚球逐步拆解，
指尖以厘米之差越过，全场远投与手臂之力相比较，
高吊球越过
篮板——并裁定体育史上最富戏剧性的一投}]
\label{pb:g12:projectile-motion:1}
篮球架的篮圈是水平圆环，离地面高 \qty{3.05}{m}；
诺拉每次都在 \qty{2.00}{m} 高处出手，以出手点为原点建立坐标轴，
水平方向 $Ox$ 指向篮筐，竖直向上为 $Oy$，篮圈中心
位于原点上方 \qty{1.05}{m} 处。除特别说明外，忽略空气阻力。

\medskip\noindent\textbf{第一部分——罚球。}诺拉在水平距离
$d = \qty{4.20}{m}$ 处投篮，\omterm{def:g12:projectile-motion:projectile}{抛射角} $\alpha = \ang{50}$。
\begin{enumerate}
  \item 出手后受哪些力？给出 $\vect a$，并注明所依据的定律。
  \item 推导 $\vect v(t)$ 的分量。
  \item 推导 $x(t)$ 与 $y(t)$。
  \item 消去 $t$，得到\omterm{def:g10:relative-motion:trajectory}{轨迹} $y(x)$。
  \item 球必须经过篮圈中心，即 $y(d) = \qty{1.05}{m}$：求
        $v_0$，并分别以 \unit{m/s} 和 \unit{km/h} 表示。
\end{enumerate}

\medskip\noindent\textbf{第二部分——越过指尖，射入篮圈。}
\begin{enumerate}[resume]
  \item 球到达篮圈需要多长时间？
  \item 求最高点的时刻与高度，并检验
        球在篮圈处已处于下降阶段。
  \item 防守者位于诺拉前方 \qty{0.90}{m} 处，伸出指尖至
        \qty{2.60}{m} 高。球越过指尖的余量为多少？
  \item 求篮圈处的 $v_x$ 与 $v_y$、该处的速率，以及速度
        在水平线以下的夹角。
  \item 为什么球必须\emph{下降}着到达篮圈？又为何要射得较陡？
\end{enumerate}

\medskip\noindent\textbf{第三部分——全场远投。}终场哨声即将
响起，诺拉距远端篮圈 \qty{26.0}{m}。
\begin{enumerate}[resume]
  \item 相对 \qty{26.0}{m} 而言，\qty{1.05}{m} 的升高很小：论证
        可采用平地模型，并由 $y(t_f) = 0$ 推导\omterm{def:g12:projectile-motion:range}{飞行时间}。
  \item 推出 $R = v_0^2\sin(2\alpha)/g$，并求\omterm{def:g12:projectile-motion:range}{射程}最大时的\omterm{def:g12:projectile-motion:projectile}{抛射角}。
  \item 求 $v_0$，使\omterm{def:g12:projectile-motion:range}{射程}达到 \qty{26.0}{m}，\omterm{def:g12:projectile-motion:projectile}{抛射角}为 $\alpha = \ang{45}$。
  \item 一名力量很强的球员可以以约 \qty{17.5}{m/s} 的速率掷出篮球。
        她的最大\omterm{def:g12:projectile-motion:range}{射程}是多少——这一投在人力上可行吗？
  \item 速率为 \qty{17.5}{m/s} 时，求能使球落在 \qty{26.0}{m} 远处的
        \emph{两个}角度，并求它们的\omterm{def:g12:projectile-motion:range}{飞行时间}。若时钟剩余 \qty{2.0}{s}，
        哪一个能赶在哨声响起前投中？
\end{enumerate}

\medskip\noindent\textbf{第四部分——越过篮板的高吊球。}诺拉在
底线身后投出高吊球：出手点距篮圈中心 \qty{6.0}{m}，
\omterm{def:g12:projectile-motion:projectile}{抛射角} $\alpha = \ang{75}$；篮板所在平面与其\omterm{def:g10:relative-motion:trajectory}{轨迹}相交于距篮圈中心 \qty{0.45}{m} 处，
位于篮圈中心之前，篮板顶边离地 \qty{3.95}{m}。
\begin{enumerate}[resume]
  \item 求使球穿过篮圈中心的 $v_0$。
  \item 计算球在篮板所在平面处的高度
        （$x = \qty{5.55}{m}$）：球能否越过篮板顶边？若能，超出多少？
  \item 求最高点高度与\omterm{def:g12:projectile-motion:range}{飞行时间}——观众为何如此喜欢这种投法？
  \item 定性说明空气阻力对这三次投篮有何改变，
        其中哪一次受其影响最大？
  \item \emph{绝杀裁定}，用两句话作答：比较全场投篮所需的速率
        与手臂所能提供的速率相比，给出
        余量，并裁定其为可能还是传说？
\end{enumerate}
\end{problem}
